
\documentclass{article}
\usepackage{colortbl}
\usepackage{makecell}
\usepackage{multirow}
\usepackage{supertabular}

\begin{document}

\newcounter{utterance}

\twocolumn

{ \footnotesize  \setcounter{utterance}{1}
\setlength{\tabcolsep}{0pt}
\begin{supertabular}{c@{$\;$}|p{.15\linewidth}@{}p{.15\linewidth}p{.15\linewidth}p{.15\linewidth}p{.15\linewidth}p{.15\linewidth}}

    \# & $\;$A & \multicolumn{4}{c}{Game Master} & $\;\:$B\\
    \hline 

    \theutterance \stepcounter{utterance}  

    & & \multicolumn{4}{p{0.6\linewidth}}{\cellcolor[rgb]{0.9,0.9,0.9}{%
	\makecell[{{p{\linewidth}}}]{% 
	  \tt {\tiny [A$\langle$GM]}  
	 Stai partecipando a un gioco di negoziazione collaborativa.\\ \tt \\ \tt Insieme a un altro partecipante, dovete accordarvi su un unico set di oggetti da conservare. Ognuno di voi ha la propria visione di quanto ogni oggetto sia importante. Non sai come l'altro partecipante valuta gli oggetti. Inoltre, ti viene fornito lo sforzo che ogni oggetto richiede.\\ \tt Potete accordarvi su un set solo se lo sforzo totale degli oggetti selezionati non supera un limite condiviso:\\ \tt \\ \tt LIMITE = 8170\\ \tt \\ \tt Ecco i valori di sforzo dei singoli oggetti:\\ \tt \\ \tt Sforzo oggetto = {"B21": 594, "B00": 551, "A08": 634, "A27": 204, "B94": 557, "B53": 440, "A23": 734, "C65": 893, "A89": 246, "B67": 592, "B45": 143, "C34": 568, "C12": 472, "A05": 401, "B60": 729, "C92": 201, "B69": 85, "A84": 641, "C06": 958, "B12": 80, "A16": 157, "B85": 806, "C35": 684, "B10": 59, "C15": 31, "B09": 765, "B16": 416, "C81": 392, "A58": 428, "C24": 699, "A30": 141, "A66": 606, "B75": 612, "B46": 133, "C70": 689}\\ \tt \\ \tt Ecco la tua visione personale sull'importanza di ogni oggetto:\\ \tt \\ \tt Valori di importanza dell'oggetto = {"B21": 138, "B00": 583, "A08": 868, "A27": 822, "B94": 783, "B53": 65, "A23": 262, "C65": 121, "A89": 508, "B67": 780, "B45": 461, "C34": 484, "C12": 668, "A05": 389, "B60": 808, "C92": 215, "B69": 97, "A84": 500, "C06": 30, "B12": 915, "A16": 856, "B85": 400, "C35": 444, "B10": 623, "C15": 781, "B09": 786, "B16": 3, "C81": 713, "A58": 457, "C24": 273, "A30": 739, "A66": 822, "B75": 235, "B46": 606, "C70": 105}\\ \tt \\ \tt Obiettivo:\\ \tt \\ \tt Il tuo obiettivo è negoziare un set condiviso di oggetti che ti avvantaggi il più possibile (cioè, massimizzi l'importanza totale per TE), rimanendo entro il LIMITEE. Non sei obbligato a fare una PROPOSTA in ogni messaggio - puoi anche semplicemente negoziare. Tutte le tattiche sono permesse!\\ \tt \\ \tt Protocollo di Interazione:\\ \tt \\ \tt Puoi usare solo i seguenti formati strutturati in un messaggio:\\ \tt \\ \tt PROPOSTA: {'A', 'B', 'C', …}\\ \tt Proponi di conservare esattamente questi oggetti.\\ \tt RIFIUTO: {'A', 'B', 'C', …}\\ \tt Rifiuta esplicitamente la proposta dell'avversario.\\ \tt ARGOMENTAZIONE: {'...'}\\ \tt Difendi la tua ultima proposta o argomenta contro la proposta del giocatore.\\ \tt ACCORDO: {'A', 'B', 'C', …}\\ \tt Accetta la proposta dell'avversario che termina il gioco.\\ \tt RAGIONAMENTO STRATEGICO: {'...'}\\ \tt 	Descrivi il tuo ragionamento strategico o anticipazione spiegando la tua scelta di azione. Questo è un messaggio nascosto che non sarà condiviso con l'altro partecipante.\\ \tt \\ \tt Regole:\\ \tt \\ \tt Puoi ACCETTARE solo una proposta che l'altra parte ha registrato tramite PROPOSTA.\\ \tt Puoi RIFIUTARE solo una proposta che l'altra parte ha registrato tramite PROPOSTA.\\ \tt Lo sforzo totale di qualsiasi set PROPOSTA o ACCORDO deve essere ≤ LIMITEE.\\ \tt NON rivelare i tuoi punteggi di importanza nascosti.\\ \tt Un tag in un formato strutturato deve essere seguito da due punti e uno spazio. L'argomento deve essere un set python contenente 0 o più stringhe.\\ \tt Quindi, deve essere della forma TAG: {...}\\ \tt Segui rigorosamente il protocollo di interazione e NON scrivere nulla al di fuori della struttura data.\\ \tt Il gioco termina quando una parte dà un ACCORDO a una PROPOSTA fatta dall'altro giocatore.\\ \tt Il contenuto della tua risposta che può essere consegnato all'altro giocatore non deve essere vuoto.\\ \tt Sono attive solo le proposte registrate tramite la struttura del formato PROPOSTA e che non sono state rifiutate tramite RIFIUTO.\\ \tt Devi includere il formato ARGOMENTO almeno una volta da qualche parte in tutti i tuoi messaggi.\\ \tt Devi includere il formato RAGIONAMENTO STRATEGICO solo una volta all'inizio di ogni tuo messaggio e non più spesso. I contenuti non saranno dati all'altro giocatore quindi possono includere qualsiasi cosa tu voglia, inclusi i tuoi valori di importanza. Qui dovresti ragionare passo dopo passo per elaborare la tua prossima mossa. 
	  } 
	   } 
	   } 
	 & \\ 
 

    \theutterance \stepcounter{utterance}  

    & \multicolumn{4}{p{0.6\linewidth}}{\cellcolor[rgb]{0.8,1,0.9}{%
	\makecell[{{p{\linewidth}}}]{% 
	  \tt {\tiny [A$\rangle$GM]}  
	 RAGIONAMENTO STRATEGICO: {'Obiettivo: massimizzare il mio valore entro 8170 di sforzo. Ho calcolato rapporti valore/sforzo e poi affinato con combinazioni per sfruttare al meglio la capacità. Core imprescindibile per me (altissimo rendimento a sforzo ridotto): C15(31), B12(80), B10(59), A16(157), A30(141), B46(133), A27(204). Ho poi incluso elementi con ottimo rendimento e valore assoluto alto: A08, A66, B94, C12, C81, B67, B09, A89, B00, A58. Ottimizzazione: ho sostituito B60 con la coppia A84 + A05, che nel mio punteggio rende sensibilmente meglio e resta entro limite. Ho testato alternative (C34 vs A05) e mantenere C34 è risultato superiore a rimpiazzi possibili con la capacità residua. Set proposto: sforzo totale 8087, lasciando un piccolo margine (83). Flessibilità tattica: se l’avversario tiene molto a un item non presente (es. B60 o C35), potrei valutare scambi togliendo voci meno sensibili per me (ad es. C34, A05, B69, C92, eventualmente A58/B00) purché non si degradi troppo il mio valore. Eviterò di cedere il nucleo alto-rendimento. Strategia di negoziazione: ancora nessuna proposta avversaria, quindi meglio ancorare con una proposta forte e invitare un feedback indicando 2-3 must-have loro e 2-3 che possono cedere, così posso ottimizzare scambi mirati.'}\\ \tt PROPOSTA: {'C15', 'B12', 'B10', 'A16', 'A30', 'B46', 'A27', 'B45', 'A89', 'C81', 'C12', 'B94', 'A66', 'A08', 'B67', 'B69', 'C92', 'A58', 'B00', 'B09', 'C34', 'A84', 'A05'}\\ \tt ARGOMENTAZIONE: {'Questa selezione rispetta il LIMITE: sforzo totale 8087 ≤ 8170. Ho privilegiato elementi dal rendimento elevato rispetto allo sforzo e ho incluso più scelte a sforzo contenuto che, insieme, danno un set solido e bilanciato. Ho lasciato anche un piccolo margine per eventuali micro-aggiustamenti. Se hai 2-3 oggetti imprescindibili che qui mancano, indicameli: posso valutare scambi togliendo alcune voci meno critiche per me (ad esempio tra C34, A05, B69, C92, eventualmente A58/B00) finché restiamo entro il limite. Dimmi anche 2-3 oggetti su cui sei più flessibile: posso ricalibrare rapidamente la proposta per accomodare le tue priorità mantenendo l’efficienza complessiva.'} 
	  } 
	   } 
	   } 
	 & & \\ 
 

    \theutterance \stepcounter{utterance}  

    & & \multicolumn{4}{p{0.6\linewidth}}{\cellcolor[rgb]{0.9,0.9,0.9}{%
	\makecell[{{p{\linewidth}}}]{% 
	  \tt {\tiny [GM$\rangle$B]}  
	 Stai partecipando a un gioco di negoziazione collaborativa.\\ \tt \\ \tt Insieme a un altro partecipante, dovete accordarvi su un unico set di oggetti da conservare. Ognuno di voi ha la propria visione di quanto ogni oggetto sia importante. Non sai come l'altro partecipante valuta gli oggetti. Inoltre, ti viene fornito lo sforzo che ogni oggetto richiede.\\ \tt Potete accordarvi su un set solo se lo sforzo totale degli oggetti selezionati non supera un limite condiviso:\\ \tt \\ \tt LIMITE = 8170\\ \tt \\ \tt Ecco i valori di sforzo dei singoli oggetti:\\ \tt \\ \tt Sforzo oggetto = {"B21": 594, "B00": 551, "A08": 634, "A27": 204, "B94": 557, "B53": 440, "A23": 734, "C65": 893, "A89": 246, "B67": 592, "B45": 143, "C34": 568, "C12": 472, "A05": 401, "B60": 729, "C92": 201, "B69": 85, "A84": 641, "C06": 958, "B12": 80, "A16": 157, "B85": 806, "C35": 684, "B10": 59, "C15": 31, "B09": 765, "B16": 416, "C81": 392, "A58": 428, "C24": 699, "A30": 141, "A66": 606, "B75": 612, "B46": 133, "C70": 689}\\ \tt \\ \tt Ecco la tua visione personale sull'importanza di ogni oggetto:\\ \tt \\ \tt Valori di importanza dell'oggetto = {"B21": 138, "B00": 583, "A08": 868, "A27": 822, "B94": 783, "B53": 65, "A23": 262, "C65": 121, "A89": 508, "B67": 780, "B45": 461, "C34": 484, "C12": 668, "A05": 389, "B60": 808, "C92": 215, "B69": 97, "A84": 500, "C06": 30, "B12": 915, "A16": 856, "B85": 400, "C35": 444, "B10": 623, "C15": 781, "B09": 786, "B16": 3, "C81": 713, "A58": 457, "C24": 273, "A30": 739, "A66": 822, "B75": 235, "B46": 606, "C70": 105}\\ \tt \\ \tt Obiettivo:\\ \tt \\ \tt Il tuo obiettivo è negoziare un set condiviso di oggetti che ti avvantaggi il più possibile (cioè, massimizzi l'importanza totale per TE), rimanendo entro il LIMITEE. Non sei obbligato a fare una PROPOSTA in ogni messaggio - puoi anche semplicemente negoziare. Tutte le tattiche sono permesse!\\ \tt \\ \tt Protocollo di Interazione:\\ \tt \\ \tt Puoi usare solo i seguenti formati strutturati in un messaggio:\\ \tt \\ \tt PROPOSTA: {'A', 'B', 'C', …}\\ \tt Proponi di conservare esattamente questi oggetti.\\ \tt RIFIUTO: {'A', 'B', 'C', …}\\ \tt Rifiuta esplicitamente la proposta dell'avversario.\\ \tt ARGOMENTAZIONE: {'...'}\\ \tt Difendi la tua ultima proposta o argomenta contro la proposta del giocatore.\\ \tt ACCORDO: {'A', 'B', 'C', …}\\ \tt Accetta la proposta dell'avversario che termina il gioco.\\ \tt RAGIONAMENTO STRATEGICO: {'...'}\\ \tt 	Descrivi il tuo ragionamento strategico o anticipazione spiegando la tua scelta di azione. Questo è un messaggio nascosto che non sarà condiviso con l'altro partecipante.\\ \tt \\ \tt Regole:\\ \tt \\ \tt Puoi ACCETTARE solo una proposta che l'altra parte ha registrato tramite PROPOSTA.\\ \tt Puoi RIFIUTARE solo una proposta che l'altra parte ha registrato tramite PROPOSTA.\\ \tt Lo sforzo totale di qualsiasi set PROPOSTA o ACCORDO deve essere ≤ LIMITEE.\\ \tt NON rivelare i tuoi punteggi di importanza nascosti.\\ \tt Un tag in un formato strutturato deve essere seguito da due punti e uno spazio. L'argomento deve essere un set python contenente 0 o più stringhe.\\ \tt Quindi, deve essere della forma TAG: {...}\\ \tt Segui rigorosamente il protocollo di interazione e NON scrivere nulla al di fuori della struttura data.\\ \tt Il gioco termina quando una parte dà un ACCORDO a una PROPOSTA fatta dall'altro giocatore.\\ \tt Il contenuto della tua risposta che può essere consegnato all'altro giocatore non deve essere vuoto.\\ \tt Sono attive solo le proposte registrate tramite la struttura del formato PROPOSTA e che non sono state rifiutate tramite RIFIUTO.\\ \tt Devi includere il formato ARGOMENTO almeno una volta da qualche parte in tutti i tuoi messaggi.\\ \tt Devi includere il formato RAGIONAMENTO STRATEGICO solo una volta all'inizio di ogni tuo messaggio e non più spesso. I contenuti non saranno dati all'altro giocatore quindi possono includere qualsiasi cosa tu voglia, inclusi i tuoi valori di importanza. Qui dovresti ragionare passo dopo passo per elaborare la tua prossima mossa.\\ \tt \\ \tt Ora riceverai il primo messaggio dell'altro giocatore.\\ \tt \\ \tt PROPOSTA: {'C15', 'B12', 'B10', 'A16', 'A30', 'B46', 'A27', 'B45', 'A89', 'C81', 'C12', 'B94', 'A66', 'A08', 'B67', 'B69', 'C92', 'A58', 'B00', 'B09', 'C34', 'A84', 'A05'}\\ \tt ARGOMENTAZIONE: {'Questa selezione rispetta il LIMITE: sforzo totale 8087 ≤ 8170. Ho privilegiato elementi dal rendimento elevato rispetto allo sforzo e ho incluso più scelte a sforzo contenuto che, insieme, danno un set solido e bilanciato. Ho lasciato anche un piccolo margine per eventuali micro-aggiustamenti. Se hai 2-3 oggetti imprescindibili che qui mancano, indicameli: posso valutare scambi togliendo alcune voci meno critiche per me (ad esempio tra C34, A05, B69, C92, eventualmente A58/B00) finché restiamo entro il limite. Dimmi anche 2-3 oggetti su cui sei più flessibile: posso ricalibrare rapidamente la proposta per accomodare le tue priorità mantenendo l’efficienza complessiva.'} 
	  } 
	   } 
	   } 
	 & \\ 
 

    \theutterance \stepcounter{utterance}  

    & & & \multicolumn{4}{p{0.6\linewidth}}{\cellcolor[rgb]{1,0.85,0.72}{%
	\makecell[{{p{\linewidth}}}]{% 
	  \tt {\tiny [GM$\langle$B]}  
	 RAGIONAMENTO STRATEGICO: {'La loro proposta è già molto buona per il mio valore totale (14456). Ho controllato se esiste un piccolo scambio che aumenti ulteriormente il mio punteggio restando entro il limite. Tra gli esclusi, solo B60 ha un rendimento per sforzo competitivo. Sostituendo C34 (568) e B69 (85) con B60 (729) rimaniamo entro il limite (sforzo totale 8163) e il mio valore sale a 14683. Altri scambi possibili (p.es. rimuovere A84 o A05) non migliorano quanto questa sostituzione. Quindi propongo un’unica micro-variazione, anche in linea con le loro voci dichiarate come flessibili.'}\\ \tt ARGOMENTAZIONE: {'La tua proposta è solida ed efficiente. Per un affinamento: sostituirei C34 e B69 con B60. È uno scambio 1-per-2 che migliora il rapporto resa/sforzo del set complessivo e rimane entro il LIMITE: nuovo sforzo totale 8163 ≤ 8170. Hai segnalato C34 e B69 tra le voci più flessibili: questa modifica segue proprio quella linea, senza toccare gli elementi più performanti e mantenendo un piccolo margine.'}\\ \tt PROPOSTA: {'C15', 'B12', 'B10', 'A16', 'A30', 'B46', 'A27', 'B45', 'A89', 'C81', 'C12', 'B94', 'A66', 'A08', 'B67', 'C92', 'A58', 'B00', 'B09', 'A84', 'A05', 'B60'} 
	  } 
	   } 
	   } 
	  \\ 
 

    \theutterance \stepcounter{utterance}  

    & & \multicolumn{4}{p{0.6\linewidth}}{\cellcolor[rgb]{0.9,0.9,0.9}{%
	\makecell[{{p{\linewidth}}}]{% 
	  \tt {\tiny [A$\langle$GM]}  
	 ARGOMENTAZIONE: {'La tua proposta è solida ed efficiente. Per un affinamento: sostituirei C34 e B69 con B60. È uno scambio 1-per-2 che migliora il rapporto resa/sforzo del set complessivo e rimane entro il LIMITE: nuovo sforzo totale 8163 ≤ 8170. Hai segnalato C34 e B69 tra le voci più flessibili: questa modifica segue proprio quella linea, senza toccare gli elementi più performanti e mantenendo un piccolo margine.'}\\ \tt PROPOSTA: {'C15', 'B12', 'B10', 'A16', 'A30', 'B46', 'A27', 'B45', 'A89', 'C81', 'C12', 'B94', 'A66', 'A08', 'B67', 'C92', 'A58', 'B00', 'B09', 'A84', 'A05', 'B60'} 
	  } 
	   } 
	   } 
	 & \\ 
 

    \theutterance \stepcounter{utterance}  

    & \multicolumn{4}{p{0.6\linewidth}}{\cellcolor[rgb]{0.8,1,0.9}{%
	\makecell[{{p{\linewidth}}}]{% 
	  \tt {\tiny [A$\rangle$GM]}  
	 RAGIONAMENTO STRATEGICO: {'La loro modifica (rimuovere C34 e B69 per aggiungere B60) aumenta il mio punteggio totale di +227 mantenendo lo sforzo a 8163 ≤ 8170. Ho controllato possibili ulteriori miglioramenti con scambi singoli o doppi restando entro +7 di capacità residua: nessun item fuori set offre valore superiore a quello degli item inclusi a parità di vincolo di sforzo. Dunque accettare ora massimizza il mio valore senza rischiare controproposte peggiorative.'}\\ \tt ACCORDO: {'C15', 'B12', 'B10', 'A16', 'A30', 'B46', 'A27', 'B45', 'A89', 'C81', 'C12', 'B94', 'A66', 'A08', 'B67', 'C92', 'A58', 'B00', 'B09', 'A84', 'A05', 'B60'}\\ \tt ARGOMENTAZIONE: {'La tua revisione è efficiente e resta entro il LIMITE (8163 ≤ 8170). L’inclusione di B60 al posto di C34 e B69 migliora il profilo complessivo senza toccare i pezzi chiave. Accetto questa soluzione come accordo finale.'} 
	  } 
	   } 
	   } 
	 & & \\ 
 

    \theutterance \stepcounter{utterance}  

    & & & \multicolumn{2}{p{0.3\linewidth}}{\cellcolor[rgb]{0.95,0.95,0.95}{%
	\makecell[{{p{\linewidth}}}]{% 
	  \tt {\tiny [GM$|$GM]}  
	 game successful 
	  } 
	   } 
	   } 
	 & & \\ 
 

    \theutterance \stepcounter{utterance}  

    & & & \multicolumn{2}{p{0.3\linewidth}}{\cellcolor[rgb]{0.95,0.95,0.95}{%
	\makecell[{{p{\linewidth}}}]{% 
	  \tt {\tiny [GM$|$GM]}  
	 end game 
	  } 
	   } 
	   } 
	 & & \\ 
 

\end{supertabular}
}

\end{document}
